\cotefile{64}
\scanpage{131}
Modules des courbes elliptiques : théorie transcendante

La donnée d'une courbe elliptique \inserted{$E$} sur $\mathbb C$ revient à celle \deleted{\unsure{à}} d'un \inserted{$\mathbb C$-}espace vectoriel de rang $1$, savoir \deleted{$\unsure{\mathfrak t_E|t_E}$}
\[
 V=\mathfrak t_E \qquad \text{(espace tangent à l'origine)} \tag*{(1)}
\]
et d'un ``réseau''
\[
 H\subset V \tag*{(2)}
\]
(sous-\deleted{groupe abélien discret}\inserted{$\mathbb Z$-module discret, libre} de rang $2$). La donnée du \deleted{groupe} $\mathbb Z$-module $\Pi$ permet de récupérer $V$ en tant qu'espace vectoriel réel \inserted{(de rang $2$)} par
\[
 V\simeq H\otimes_{\mathbb Z}\mathbb R\mathrel{\underset{\unsure{\circ}}{\overset{\mathrm{def}}{=}}}H_{\mathbb R} \tag*{(3)}
\]
et $E$ en tant que groupe de Lie réel, par
\[
 E\simeq V/H\simeq H_{\mathbb R}/H\,, \tag*{(4)}
\]
tandis que la structure complexe de $E$ est déterminée par la donnée d'une structure d'espace vectoriel complexe sur $V$ par passage aux quotients.

Étant donné un espace vectoriel réel $V$\inserted{\deleted{de rang \unsure{$n$}}}, la donnée d'une structure complexe sur $V$ (i.e. d'un \unsure{autom|endom} $u$ de $V$ de carré $-\mathrm{id}_V$) équivaut aussi à celle d'un \inserted{$\mathbb C$-}sous-espace vectoriel $L$ de \deleted{$\unsure{V_{\mathbb C}}$}
\[
 V_{\mathbb C}=V\otimes_{\mathbb R}\mathbb C \tag*{(5)}
\]
tel que la projection canonique
\[
 \begin{aligned}
 V_{\mathbb C}&\longrightarrow V=V_{\mathbb R}\\
 z&\longmapsto \tfrac12(z+\bar z)=\Re z
 \end{aligned} \tag*{(6)}
\]
\deleted{soit} \inserted{induise} un isomorphisme d'espaces vectoriels réels.
\[
 \text{\deleted{Pour un}}\quad\varphi\colon L\xrightarrow{\sim}V\quad(\mathbb R\text{-iso}) \tag*{(7)}
\]
% archive sheet number: 130
\scanpage{132}
A un tel $L$, on associe la structure complexe de \deleted{$V$ déduite \inserted{\deleted{\unsure{conjuguée}}} de celle \inserted{\deleted{\unsure{déduite}}} de $L$ par transport de} \deleted{\unsure{conjuguée} de celle déduite de} $V$ qui est conjuguée de celle déduite de la str. complexe de $L$ par transp de structure via (7), i.e. on pose
\[
 ux=\varphi(-i\varphi^{-1}(x))=-\varphi(i\varphi(x)) \tag*{(8)}
\]
$u$ étant la multiplication par $i$ pour la structure complexe envisagée sur $V$.

% slanted
\margin{cette convention est imposée par la situation de ``modules'' des courbes elliptiques et du fait qu'on veut que les paramètres qu'on \unsure{(va)} y introduit soient \uline{holomorphes}, et non antiholomorphes \ldots}

Inversement, si $V$ est muni d'une structure complexe, \deleted{l'applic} faisant de $V$ un espace vectoriel complexe $V_c$, alors l'application \inserted{identique} $\mathbb R$-linéaire $V\longrightarrow V_c$ se prolonge, par la propriété universelle de $V_{\mathbb C}$, en une application $\mathbb C$-linéaire
\[
 \psi\colon V_{\mathbb C}\longrightarrow V_c \tag*{(9)}
\]
de sorte que
\[
 \left\{\begin{aligned}\psi(ix)&=ux&&\forall x\in V=V_{\mathbb R},\\ \psi(x)&=x\end{aligned}\right. \tag*{(10)}
\]
donc
\[
 \psi(x+iy)=x+uy \tag*{(11)}
\]
\deleted{Considérons} \inserted{\deleted{\unsure{Posons|Prenons}}} Considérons
\[
 \begin{aligned}
 L^{\unsure{\prime|1}}=\mathrm{Ker}\,\psi
 &=\{x+iy\mid x,y\in V_{\mathbb R},\ x+uy=0\quad\text{i.e. }y=\deleted{\unsure{-}}ux\}\\
 &=\{\deleted{\unsure{uy+iy}}\inserted{x+iu(x)}\mid x\in V_{\mathbb R}\}
 \end{aligned}\tag*{(12)}
\]
et considérons l'application correspondante (6)
\[
 \varphi\colon\underset{\substack{\Vert\\x+iu(x)}}{z}\longmapsto\tfrac12(z+\bar z)=\deleted{\unsure{x+uy}}\inserted{\Re z=x}\colon L\longrightarrow V
\]
On voit alors que cette application est bijective, et \uline{antilinéaire}, \inserted{donc \deleted{la} correspond à la structure complexe} \deleted{\unsure{Prenant une|Pour un} \unsure{pour $V$ de dim. finie} sur $\mathbb R$, $L=L'=\overline{\mathrm{Ker}\,\psi}$ \unsure{$\subset V_{\mathbb C}$}, \unsure{$\mathrm{dim}_{\mathbb R}=2|\mathrm{rang}_{\mathbb R}=2$},}
% slanted
\margin{mise sur $V$. alors $L\longrightarrow V_c$ est \uline{\unsure{linéaire|antilinéaire}}}

\[
 \tag*{(13)}
 \left[\begin{minipage}{0.84\linewidth}
 On peut dire aussi, que la donnée d'une structure complexe sur $V$ équivaut à

 $1^\circ$) La donnée d'un \deleted{\unsure{opé}} sous-groupe compact \deleted{maximal} $U$ \deleted{de $\mathrm{Gl}=$} $\inserted{G=}\mathrm{Aut}_{\mathbb R}(W)$, \deleted{i.e. d'un}
 \end{minipage}\right.
\]
% archive sheet number: 131
\scanpage{133}
\deleted{\unsure{sous-groupe}} isomorphe à $\mathbb U=\{z\in\mathbb C\mid |z|=1\}$ tel que l'action sur $V_{\mathbb C}$ ait exactement 2 sous-espaces propres,
% slanted
\margin{(c'est automatique si $\mathrm{rg}_{\mathbb R}V=2$)}

$2^\circ$) D'un des deux éléments d'ordre $4$ de $\mathbb U$, ou ce qui revient au $=$, d'une orientation de $\mathbb U$, ou encore \inserted{(si $\mathrm{rg}_{\mathbb R}V=2$, plus général\textsuperscript{t} si $\mathrm{rg}_{\mathbb R}V=2n$, $n$ impair)} d'une orientation de $V$
% slanted
\margin{ou encore d'un \inserted{des deux} isomorphismes $\mathbb U\xrightarrow{\sim}U$, ou $U\xrightarrow{\sim}\mathbb U$}

-- l'élément d'ordre $4$ de $\mathbb U$ associé à une orientation $\omega$ de \inserted{$U$ (ou de} $V$) sera noté $u_\omega$, c'est donc la rotation de $+\pi/2$ pour l'orientation \inserted{$\omega$} de $V$ et \unsure{leur|la|une} structure \deleted{complexe}\inserted{euclidienne} sur $V$ choisie invariante sous $U$ (\uline{NB} \inserted{si $\mathrm{rg}_{\mathbb R}V=2$,} une telle \deleted{struct}\inserted{métrique} est unique à scalaire multiplicatif près).

Alors $V_{\mathbb C}$ sous l'action de $U$ se décompose en deux sous-espaces propres
\[
 V_{\mathbb C}=\bigoplus_{\omega\in\underline\omega}\deleted{\unsure{\omega}}L_\omega \tag*{(\unsure{14|13})}
\]
où pour $\omega\in\underline\omega(U)=\underline\omega$, $L_\omega$ est caractérisé par le fait que l'isomorphisme
\[
 \varphi_\omega\colon U\xrightarrow{\sim}\mathbb U\qquad g(z)=\varphi_\omega(g)\cdot z\quad\text{pour }z\in L_\omega \tag*{(15)}
\]
défini par l'action de $U$ sur $L_\omega$, corresponde à l'orientation $\omega$. La donnée d'une structure complexe sur $V$ \inserted{compatible avec $U$} équivaut alors à celle d'un des deux isomorphismes $\varphi_\omega$ (permettant de faire opérer $\mathbb U$ sur $V$, donc aussi $\mathbb R^{\times+}\times\mathbb U\simeq\mathbb C^\times$, \deleted{\unsure{donc|dans}} $\mathbb C$ \ldots). Ceci dit, pour \deleted{$\unsure{x\in\mathbb C|z\in\mathbb C}$} $z\in L_\omega$, on aura
\[
 \begin{aligned}
 \Re (\varphi_\omega(g)z)=\Re (gz)=g\Re (z)&=\text{produit de}\\
 &\Re (z)\ \text{par }\varphi_\omega(g)\ \text{pour la structure complexe}\\[-2pt]
 &\text{de }V\text{ associée à }\varphi_\omega
 \end{aligned}
\]
% archive sheet number: 132
\scanpage{134}
i.e. \(\unsure{1|\vert}\)
\[
 \begin{aligned} L_\omega&\simeq V\\ z&\longmapsto\Re z\end{aligned}
\]
est un isomorphisme d'espaces vectoriels complexes, quand on munit $V$ de la structure complexe $c_\omega$.\(\big]\)
% slanted
\margin{Mais il y a lieu d'associer à cette structure complexe sur $V$ \unsure{non pas|définie par|non par} $L_\omega$, mais $L_{\omega'}$, cf plus bas \ldots}

Dans le cas qui nous occupe, $\mathrm{rg}_{\mathbb R}V=2$ les \deleted{droites} sous-espaces vectoriels envisagés de $V_{\mathbb C}$ sont de \deleted{\unsure{rg}} rang $1$, \inserted{corang $1$,} et peuvent s'interpréter comme les points de
\[
 \Sigma=\Sigma_V=\mathbb P(\mathbb C),\quad\text{où}\quad P=\mathbb P(V) \tag*{(16)}
\]
est \deleted{l'}le \deleted{\unsure{esp}} schéma droite projective réelle associé à $V$. Pour que $\begin{aligned}L&\longrightarrow V\\ z&\longrightarrow\Re z\end{aligned}$ soit un isomorphisme, \deleted{i.e.} pour que $L$ corresponde à une structure complexe \deleted{de} sur $V$, il f. et s. que \deleted{\unsure{le} le point $z$ correspondant de $\Sigma$ soit} \inserted{$L\cap V=\{0\}$ i.e. $L$} n'est pas complexifié d'une \deleted{\unsure{ne provi|espèce|espace}} droite $L_0\subset V_{\deleted{\unsure{\mathbb R}}}$, i.e. que le point $z\in\Sigma$ qui correspond à $L$ \inserted{ne} soit pas dans \deleted{$\Sigma(\mathbb R)$}.
\[
 \Sigma_{\mathbb R}=P(\mathbb R)
\]
ou encore qu'on ait
\[
 z\in\Sigma\smallsetminus\Sigma_{\mathbb R}, \tag*{(17)}
\]

\deleted{\unsure{En résu|En résumé}} Ainsi, les structures complexes
% archive sheet number: 133
\scanpage{135}
sur $V$ sont en correspondance biunivoque avec l'ens $\Sigma\smallsetminus\Sigma_{\mathbb R}$ des points non réels de $\Sigma$.

Notons que l'ens $\Sigma\smallsetminus\Sigma_{\mathbb R}$ a deux composantes connexes \inserted{$\Sigma_\omega^*$}, qui correspondent (par Stokes) aux deux orientations de $\Sigma_{\mathbb R}$, ou encore aux deux orientations de $V$.

La composante connexe d'un $z\in\Sigma\smallsetminus\Sigma_{\mathbb R}$ correspond à l'orientation de $V$ définie de façon évidente par la structure complexe envisagée sur $V$ (définie par $z$).
% slanted
\margin{C'est faux, cf calcul plus bas}

Revenant à la description transcendante des courbes elliptiques, on voit que \inserted{la donnée d'}une telle courbe $E$ revient à la donnée

a) D'un $\mathbb Z$-module libre de rang $2$ $H$

b) D'une \deleted{droite} orientation \inserted{$\omega$} de $V=\Pi\otimes_{\mathbb Z}\mathbb R$, ou ce qui revient au $=$, d'une des deux bases de $\bigwedge_{\mathbb Z}^2H$ (i.e. d'une restriction du groupe structural $\mathrm{Gl}(2,\mathbb Z)$ de $H$ : $\mathrm{Sl}(2,\mathbb Z)$), \inserted{\deleted{\unsure{soit $\omega$|soit}} notée aussi $\omega$}
% slanted
\margin{On a $H=\pi_1(E,\unsure{0_E|e_E})$ $V=H\otimes_{\mathbb Z}\mathbb R\simeq\mathfrak t_E$}

c) D'un point $z\in\Sigma_\omega(V)$,

Une \uline{rigidification de Torelli} de la courbe elliptique $E$ est la donnée d'un isomorphisme
% archive sheet number: 134
\scanpage{136}
\[
 \mathbb Z^2\longrightarrow H \tag*{(18)}
\]
compatible avec les orientations, i.e. tel que l'on ait
\[
 z_1\wedge z_2=\omega \tag*{(19)}
\]
où $z_1,z_2$ sont les images par (18) des éléments de la base canonique de $\mathbb Z^2$. Cet isomorphisme (18) définit un isomorphisme correspondant
\[
 \mathbb R^2\simeq V \tag*{(20)}
\]
d'où des isomorphismes
\[
 \left\{\begin{aligned}
 \mathbb P^1_{\mathbb R}&\simeq\mathbb P^1(V),\quad\mathbb P^1(\mathbb C)\simeq\Sigma_V\\
 \mathscr D&\simeq\Sigma_{V,\omega}^*
 \end{aligned}\right. \tag*{(21)}
\]
où $\mathscr D$ est le demi-plan de Poincaré, i.e. \deleted{l'ens} la variété complexe des $z\in\mathbb C$ tels que $\deleted{\Im z}\Im z>0$
\[
 \mathscr D=\{z\in\mathbb C\mid\Im z>0\}, \tag*{(22)}
\]
identifié à $\Sigma_{\mathbb R^2,\omega_0}^*$ en identifiant $z\in\mathscr D$ à la droite
\[
 L_z\deleted{\unsure{\mathbb P_z}}=\mathbb C\cdot\underbrace{(z,1)}_{ze_1+e_2}. \tag*{(23)}
\]
% slanted
\margin{$(e_1,e_2)$ base canonique de $\mathbb R^2$. Si on associait à $z\in\mathbb C$ la droite $\mathbb C(1,z)$, alors \unsure{à|ce serait au|correspond} $\Sigma^*_{\mathbb R^2,\omega_0}$ \unsure{correspondrait|correspondant|donnerait} le \deleted{\unsure{demi}} demi-plan $\mathscr D^-$ au lieu de $\mathscr D=\mathscr D^+$}
\margin{Ainsi, les classes d'isomorphie de courbes elliptiques avec \unsure{rigidif-|rigidité} de Torelli sont en corr. biunivoque avec l'ens des pts de $\mathscr D$.}

Sur l'ens des classes d'isomorphie des courbes elliptiques $E=(\Pi,\omega,L\subset H\otimes_{\deleted{\unsure{\mathbb R}}\mathbb Z}\mathbb C)$, avec rigidification de Torelli $\mathbb Z^2\simeq H$, le groupe $\mathrm{Sl}(2,\mathbb Z)$ opère à droite de façon évidente, par composition \inserted{à droite} de (18) avec les automorphismes de $\mathbb Z^2$. \deleted{\unsure{Quand on|Quant}}\inserted{On va}
% archive sheet number: 135
\scanpage{137}
expliciter \deleted{\unsure{l}}cette action en tant qu'action sur $\mathscr D$ (défini par (22)).

Si une structure complexe sur $\mathbb R^2$ est définie par l'élément $z\in\mathscr D$, $z=\deleted{\unsure{\Re}}x+iy$, alors l'application \inserted{$\mathbb C$-linéaire} $\varphi\colon L_z=\mathbb C\cdot(z,1)\longrightarrow V$ s'identifie à
\[
 \underset{\substack{\Vert\\a+ib}}{\lambda}\longmapsto
 \Re\bigl(\deleted{\unsure{\lambda(z,1)}}\underbrace{\lambda(ze_1+e_2)}_{[(ax-by)+i(ay+bx)]e_1+(a+ib)e_2}\bigr)=(ax-by)e_1+ae_2
\]
les images inverses $\lambda_1,\lambda_2$ de $e_1,e_2$ par celle-ci sont données par
\[
 \lambda_1=-\frac{i}{y}\qquad\lambda_2=1+i\frac{x}{y}
\]
d'où
\[
 \lambda_2/\lambda_1=i(y+ix)=-x+iy
\]
donc on a
\[
 \underbrace{e_2/e_1}_{\substack{\text{quotient dans }\mathbb C^*\\
 \text{au sens de la}\\\text{structure compl. de }\mathbb R^2\\
 \text{définie par }z=x+iy\in\mathscr D}}
 =(-x+iy)^{\overline{\phantom{x}}}=-x-iy=\underset{\in\mathscr D^-}{-z}
\]

Donc si une structure de Torelli sur $E$ est définie par la \deleted{\unsure{élément|éléments}}\inserted{base} $(z_1,z_2)$ de $\Pi\subset\mathfrak t_E$, alors l'élément $z\in\mathscr D$ correspondant est donné par
\[
 \deleted{z_2/z_1=-x+iy=-\bar z\quad\text{où }z=x+iy.}
 \quad\inserted{z=-(z_2/z_1)} \tag*{(24)}
\]

Soit alors
\[
 g\in\mathrm{Sl}(2,\mathbb Z)\qquad
 g=\begin{pmatrix}a&b\\c&d\end{pmatrix}\qquad
 \begin{aligned}a,b,c,d&\in\mathbb Z\\ad-bc&=1\end{aligned}
\]
et soit \deleted{\unsure{un|une} \unsure{$u$|une}}
\[
 u\colon\mathbb Z^2\xrightarrow{\sim}\Pi
\]
une structure de Torelli, correspondant à
% archive sheet number: 136
\scanpage{138}
\[
 z_1=u(e_1),\ z_2=u(e_2)\in H\subset V
\]
alors la structure de Torelli $u\circ g$ correspond à $z'_1,z'_2\in\Pi$ donnée par
\[
 \begin{aligned}
 z'_1&=ug(e_1)=u(ae_1+ce_2)=az_1+cz_2\\
 z'_2&=ug(e_2)=u(be_1+de_2)=bz_1+dz_2,
 \end{aligned}
\]
qui correspond donc à un élément
\[
 z'=\unsure{\alpha'|x'}+iy'\in\mathscr D'
\]
donné par
\[
 z'\deleted{\underbrace{\unsure{=x'+iy'|=-x'+iy'}}_{\unsure{\tilde z'|z'}\; z'=-\bar z'}}
 =-\frac{z'_2}{z'_1}=-\frac{bz_1+dz_2}{az_1+cz_2}
 =-\frac{d\frac{z_2}{z_1}+b}{c\frac{z_2}{z_1}+a}
 =\frac{d\unsure{z|\bar z}\unsure{-|+}b}{-c\unsure{z|\bar z}+a}
 \tag*{\deleted{\unsure{(24)}} (25)}
\]
% slanted
\margin{\deleted{d'où $-z'=\frac{-dz+b}{-cz+a}$ \quad (25) $z'=\frac{dz-b}{-cz+a}$}}
où \deleted{$\unsure{\tilde z=-x+iy=-\bar z}$.} Si on \deleted{\unsure{reg}} préfère expliciter \deleted{l'}l'action à gauche $u\longmapsto u\circ g^{-1}$ de $\mathrm{Sl}(2,\mathbb Z)$ sur l'ens. des structures de Torelli mod iso, on note que
\[
 g^{-1}=\begin{pmatrix}d&-b\\-c&a\end{pmatrix}\qquad \text{\deleted{d'où,}}
\]
d'où maintenant, pour $z'=x'+iy'$ correspondant à $u\circ g^{-1}$
\[
 \deleted{\tilde z'=\bar z'=\frac{a\tilde z-b}{-c\tilde z+d}
 =\frac{-a\bar z-b}{-c\bar z-d}
 =\frac{\overline{a\bar z+b}}{\overline{c\bar z+d}}} \tag*{\deleted{\unsure{(25bis)}}}
\]
\[
 z'=\frac{az+b}{cz+d} \tag*{(25bis)}
\]
\uline{victoire}, c'est la formule habituelle, mais on fera attention que le nb\textsuperscript{re} complexe $z\in\mathscr D$ attaché à la structure de Torelli définie par \deleted{$z_1,z_2\in\mathfrak t_E$}\inserted{les deux} éléments $z_1,z_2$ de la droite complexe \deleted{$\unsure{\mathfrak t_{\mathbb R}=V}$} $\mathfrak t_E$ est donné, non
% archive sheet number: 137
\scanpage{139}
\deleted{\unsure{comme} $z_2/z_1$, \unsure{mais}} \inserted{par la formule qu'on} pourrait penser $z=z_2/z_1$, \deleted{\unsure{à}} mais par (24\deleted{\unsure{bis}}) i.e. par
\[
 z=-(z_2/z_1)^{\unsure{\overline{\phantom z}|\deleted{-}|-}}=-\unsure{\bar z_2|z_2}/\unsure{\bar z_1|z_1}\ \deleted{\in\mathscr D}.
\]
Si \deleted{\unsure{donc}} on voulait \unsure{répérer|repérer} la structure de Torelli non pas par le point $z\in\mathscr D$ choisi, mais par le nb\textsuperscript{re} complexe $\tau=z_2/z_1$ lui même (symétrique de $z$ p.r. à l'axe imaginaire), on trouverait une formule de transformation d'opération de $\mathrm{Sl}(2,\mathbb Z)$ sur $\mathscr D$ qui n'est \uline{pas} la formule habituelle.

% slanted
\margin{Donc cela montre que la condition $z\in\mathscr D^+$ signifie que $(z_1,z_2)$ est, non une \inserted{base} directe, mais une \inserted{base} \uline{inverse} de $\mathfrak t_E$. Si on \deleted{voulait}\inserted{veut} faire correspondre à $z\in\mathscr D^+$ une structure de Torelli (directe), il y a lieu de changer la convention (23) en $L_z=\mathbb C.(1,z)$. qui donne une \deleted{\unsure{orientation}}\inserted{structure complexe sur $\mathbb R^2$} \deleted{\unsure{de}} de $\mathbb R^2$ compatible avec l'orientation (et $L_z\in\Sigma^-_{\mathbb R^2,\omega_0}$, \unsure{quitte à altérer $\omega$|quitte à altérer} bien sûr). \uline{$z=-(z_2/z_1)^{-1}$}$=-\frac{z_1}{z_2}$ \deleted{\unsure{(c'est un peu}} court -- il faudrait \unsure{cf|il}, vérifier que la famille en question est de façon bien analytique et naturelle \ldots}

On trouve ici le

\uline{Théorème} \quad $\mathscr D$ \inserted{$=\{z\in\mathbb C\mid\Im z>0\}$,} muni de la famille \inserted{analytique} de courbes elliptiques \deleted{$\unsure{\mathcal E}$}
\[
 E_z=\mathbb C/(\mathbb Z.1+\mathbb Z.(-\bar z))
\]
est un revêtement universel de $M_{1,1}^{\mathrm{an}}$, correspondant à la classification des courbes elliptiques avec rigidification de Torelli ; le groupe des $M_{1,1}^{\mathrm{an}}$-automorphismes de $\mathscr D$ (opérant à droite) s'identifie au groupe $\mathrm{Sl}(2,\mathbb Z)$, opérant sur $\mathscr D$ par
\[
 gz=\frac{az+b}{cz+d}\qquad\text{pour}\quad g=\begin{pmatrix}a&b\\c&d\end{pmatrix}\in\mathrm{Sl}(2,\mathbb Z),
\]
qui correspond aussi à l'opération à g. naturelle
% archive sheet number: 138
\scanpage{140}
de $\unsure{\mathrm{Gl}|\mathrm{Sl}}(2,\mathbb Z)$ sur les structures de Torelli, par $g\cdot u=u\circ g^{-1}$ ($u\colon\mathbb Z^2\xrightarrow{\sim}\Pi_E$).

\uline{Corollaire 1} On a un iso. canonique
\[
 \mathscr T_{1,1}\simeq\mathrm{Sl}(2,\mathbb Z) \tag*{(26)}
\]
(où $\mathscr T_{11}=\pi_1(M_{11})$ est calculé relativement au revêtement universel choisi $\mathscr D$ de $M_{1,1}$).
\[
 \text{On a un iso. can}\quad M_{11}\simeq(\deleted{\unsure{\Sigma}}\mathscr D_{\deleted{\unsure{1}}},\mathrm{Sl}(2,\mathbb Z) \tag*{(27)}
\]

Utilisant les grands théorèmes \unsure{marteaux-pilons}, on en déduit

\uline{Corollaire 2} Soit $(\deleted{\unsure{\Sigma|E}}E,a)$ une surface \inserted{ponctuée} orientée compacte connexe de genre $1$, et soit $\mathscr T$ le groupe des classes d'isotopie des automorphismes de $(E,a)$. \deleted{Alors Considérons} \inserted{Alors} l'homomorphisme naturel
\[
 \mathscr T\longrightarrow\mathrm{Aut}_{\mathbb Z}^{+}(H)\qquad\text{où}\quad H=\pi_1(E,a) \tag*{(27)}
\]
est un isomorphisme. \unsure{\S|\deleted{§}}

\uline{Corollaire 3} Soit $\Pi$ un groupe à lacets de type \deleted{\unsure{$g$}} $(1,1)$, $H=\Pi_{\mathrm{ab}}$ (qui est un $\mathbb Z$-module libre de rang $1$), $\mathscr T=\mathrm{Autext}_{\unsure{\mathrm{lac}|\mathrm{loc}}}^+(\Pi)$, \deleted{considérons}\inserted{alors,} l'hom
\[
 \mathscr T\longrightarrow\mathrm{Aut}_{\mathbb Z}^+(H) \tag*{(28)}
\]
est iso.
% slanted
\margin{Il faudrait expliciter l'action \inserted{extérieure} de $\mathrm{Aut}_{\mathbb Z}^+(H)$ sur $\Pi$, \unsure{resp.|puis} l'action $\unsure{-|\text{de}}\ \mathrm{Aut}_{\mathbb Z}^+(H)$ \ldots}
% archive sheet number: 139
\scanpage{141}

\margin{\deleted{$\unsure{P\mathscr D}$}}

Nous pouvons expliciter le fibré tangent le long des \deleted{fibres} \inserted{la section nulle} de la courbe elliptique universelle rigidifiée à la Torelli, comme étant le fibré formé des $\mathfrak{t}_E$, qui du point de vue fibré vectoriel réel s'identifie au \deleted{fibré} $\mathrm{Sl}(2,\mathbb{Z})$-fibré sur $\mathscr D$ défini par l'opération \unsure{standart|standard} de $\mathrm{Sl},\mathbb{Z})$ sur $\mathbb{R}^2$. Le fibré principal associé \inserted{$P\mathscr D$} s'obtient donc (pour sa seule structure topologique) en remplaçant ds cet énoncé $\mathbb{R}^{2}$ par $(\mathbb{R}^2)^*$. Pour en trouver un revêtement universel, il suffit de choisir un revêtement universel de $(\mathbb{R}^{\unsure{2|1}})^*\simeq\mathbb{C}^*$, p. ex. lé revêtement par l'exponentielle complexe
\[
z\mapsto\exp(2i\pi z):\quad\mathbb{C}\longrightarrow\mathbb{C}^*
\]

% archive sheet number: 140
\scanpage{143}

Une courbe elliptique $E$ sur $\mathbb{C}$ est déterminée par
\[
\left\{\begin{aligned}
&\text{a) ``réseau'' }H,\ \deleted{\unsure{\text{le}}}\ \mathbb{Z}\text{-module libre de rang }2\quad\deleted{\unsure{(H=1}}\\
&(H=\pi_1(E,0_E)\simeq H_1(E,\mathbb{Z})=H^1(E,\mathbb{Z})^\vee)\\
&\text{b) Une structure complexe sur }V=H\otimes_{\mathbb{Z}}\mathbb{R}\ \deleted{\unsure{\simeq V|=V}}\ (\simeq\mathfrak{t}_E,\\
&\qquad\text{espace tangent à l'origine de }E).
\end{aligned}\right.
\tag*{(1)}
\]
La donnée a) permet de reconstituer $E$ comme groupe de Lie réel par
\[
E\simeq H\otimes_{\mathbb{Z}}\mathbb{R}\,/\,H
\tag*{(2)}
\]
et la donnée b) \unsure{redonne|redonnent} la structure complexe sur $E$, par passage au quotient.

Quand on est sur une base $S$, espace $\mathbb{C}$-analytique, la donnée a) correspond à un système local. $\mathscr{H}$ sur $S$ -- faisceau en $\mathbb{Z}$-modules loc libre de rang $2$ sur $S$, d'où un fibré vectoriel $\mathbb{R}$-analytique $\mathscr{H}\otimes_{\mathbb{Z}}\mathbb{R}$ sur $S$. La structure b) correspond à une structure complexe \deleted{sur} de ce dernier, en tant que fibré vectoriel \inserted{complexe} de rang $1$, de façon que les sections locales de $\mathscr{H}$ deviennent des sections analytiques complexes. Si \deleted{\unsure{$a$|$e$}} $z_1,z_2$ sont deux sections de $\mathscr{H}$ formant

% archive sheet number: 142
\scanpage{144}

une base de $\mathscr{H}\otimes_{\mathbb{Z}}\mathbb{R}$ (\inserted{fibre par fibre} sur $\mathbb{R}$), la donnée d'une telle structure complexe équivaut à la donnée de
\[
\tau=z_2/z_1
\tag*{(3)}
\]
en tant que section du faisceau structural $\mathcal{O}_S$. (\uline{NB} \deleted{il faut,} pour une structure complexe donnée sur $\mathscr{H}\otimes_{\mathbb{Z}}\mathbb{R}$, $z_1$ et $z_2$ sont deux sections\inserted{-bases} du faisceau inversible obtenu, et $\tau$ est la section de $\underline{\mathcal{O}}_S$ qui fait passer de l'une à l'autre :
\[
z_2=\tau z_1\,.
\tag*{(3 bis)}
\]
% slanted
\margin{Noter que $\forall s\in S$, $\tau(s)\in\mathbb{C}\smallsetminus\mathbb{R}$}

Une structure \uline{de Torelli} sur une \deleted{\unsure{fam}} courbe elliptique relative $E$ sur $\unsure{S_a|S{,}}$ (correspondant à $\mathscr{H}$ sur $S$ $+$ structure complexe sur $\mathscr{H}\otimes_{\mathbb{Z}}\mathbb{R}$) est la donnée d'un isomorphisme
\[
t:\mathbb{Z}_S^2\xrightarrow{\sim}\mathscr{H}\,,
\tag*{(4)}
\]
\deleted{\unsure{ou bien|i.e.}} \inserted{\deleted{\unsure{donc de|deux d}} ou encore} de deux sections $z_1,z_2$ de $\mathscr{H}$, \deleted{\unsure{satisfais|distinguées}} formant une base de $\mathscr{H}$ en \unsure{tte} fibre. La catégorie des courbes elliptiques relatives sur $S$, avec structure de Torelli, est \deleted{\unsure{équivalente à}} \inserted{``rigide'' \deleted{et}} (pas d'autom. $\ne\mathrm{id}$ d'un objet) et l'ens des

% archive sheet number: 143
